\section*{第 \ref{ch:b3:complex-spectra-magnetism} 章——配合物的电子光谱与磁性}

\begin{solution}{exo:b3:complex-spectra-magnetism:1}
D: E$_g$ + T$_{2g}$ ($2 + 3 = 5$); F: A$_{2g}$ + T$_{1g}$ + T$_{2g}$ ($1 + 3 + 3 = 7$); G: A$_{1g}$ + E$_g$ + T$_{1g}$ + T$_{2g}$
($1 + 2 + 3 + 3 = 9$).
\end{solution}

\begin{solution}{exo:b3:complex-spectra-magnetism:2}
$d^1$ $^2$D $\to$ $^2$T$_{2g}$; $d^2$ $^3$F $\to$ $^3$T$_{1g}$; $d^3$ $^4$F $\to$ $^4$A$_{2g}$; $d^4$ $^5$D $\to$ $^5$E$_g$; $d^5$ $^6$S $\to$ $^6$A$_{1g}$; $d^6$ $^5$D
$\to$ $^5$T$_{2g}$; $d^7$ $^4$F $\to$ $^4$T$_{1g}$; $d^8$ $^3$F $\to$ $^3$A$_{2g}$; $d^9$ $^2$D $\to$ $^2$E$_g$.
\end{solution}

\begin{solution}{exo:b3:complex-spectra-magnetism:3}
高自旋 \ce{Mn^{2+}}（$n = 5$）$5.92\,\mu_B$，\ce{Fe^{2+}}（4）$4.90\,\mu_B$，\ce{Co^{2+}}（3）$3.87\,\mu_B$；低自旋 \ce{Fe^{2+}} 和
\ce{Co^{3+}}（$t_{2g}^6$）：0，抗磁性。
\end{solution}

\begin{solution}{exo:b3:complex-spectra-magnetism:4}
$\ce{[Mn(H2O)6]^{2+}}$（自旋和宇称禁阻，几乎无色）$<$ $\ce{[Co(H2O)6]^{2+}}$
（宇称禁阻，\omterm{def:b3:complex-spectra-magnetism:vibronic-coupling}{电子振动耦合}）$<$ $\ce{[CoCl4]^{2-}}$（无对称中心）$<$ \ce{MnO4-}（允许的
电荷转移）。
\end{solution}

\begin{solution}{exo:b3:complex-spectra-magnetism:5}
8500：$^3$A$_{2g}\to{}^3$T$_{2g}$，故 $\Delta_{\mathrm o} = \qty{8500}{cm^{-1}}$；14100：$^3$T$_{1g}$(F)；24900：$^3$T$_{1g}$(P)。
$15B = 39\,000 - 3 \times 8500$, $B = \qty{900}{cm^{-1}}$, $\beta = 900/1056 = 0.85$.
\end{solution}

\begin{solution}{exo:b3:complex-spectra-magnetism:6}
$\Delta_{\mathrm o} = \qty{17400}{cm^{-1}}$；解方程 $7.5B + 1.5\Delta_{\mathrm o} - \frac12\sqrt{225B^2 - 18B\Delta_{\mathrm o} + \Delta_{\mathrm o}^2} = 24\,600$，得 $B = \qty{729}{cm^{-1}}$，
$\beta = 0.79$。$\nu_3 = \qty{38500}{cm^{-1}}$（\qty{260}{nm}），位于紫外区。
\end{solution}

\begin{solution}{exo:b3:complex-spectra-magnetism:7}
与水分子相比，$d$ 电子更多地离域到宝石中的氧离子上，
即此处 Cr--O 键的共价性更强。
\end{solution}

\begin{solution}{exo:b3:complex-spectra-magnetism:8}
强（$e_g$ 不均匀填充）：$d^9$、高自旋 $d^4$、低自旋 $d^7$。无：$d^3$（$t_{2g}^3$，A 谱项）和
低自旋 $d^6$（$t_{2g}^6$）。弱：高自旋 $d^6$（$t_{2g}^4e_g^2$，T 谱项）。
\end{solution}

\begin{solution}{exo:b3:complex-spectra-magnetism:9}
八面体高自旋 $d^7$：基态谱项 $^4$T$_{1g}$，为带\omterm{def:b3:complex-spectra-magnetism:moment}{轨道贡献}的 T 谱项，
磁矩远大于\omterm{def:b3:complex-spectra-magnetism:moment}{唯自旋磁矩} $3.87\,\mu_B$。四面体 $d^7$（$\equiv$ 八面体 $d^3$）：基态
谱项 $^4$A$_2$，轨道磁矩被猝灭，接近\omterm{def:b3:complex-spectra-magnetism:moment}{唯自旋磁矩}（
由于与激发态 T 谱项混合，略高于唯自旋值）。
\end{solution}

\begin{solution}{exo:b3:complex-spectra-magnetism:10}
$\mu_{\mathrm{eff}} = 2.828\sqrt{1.87} = 3.87\,\mu_B$：三个未成对电子（$S = \frac32$）。
\end{solution}

\begin{solution}{exo:b3:complex-spectra-magnetism:11}
$\Delta\chi_v = 3 \times 25.0/\num{400e6} = \num{1.88e-7}$；$\chi_{\mathrm m} = \num{1.88e-7}/\qty{10.0}{mol.m^{-3}} = \qty{1.88e-8}{m^3.mol^{-1}}$，即
\qty{1.49e-3}{cm^3.mol^{-1}}；$\chi_{\mathrm m}T = 0.445$，$\mu_{\mathrm{eff}} = 1.89\,\mu_B$：一个未成对电子。
\end{solution}

\begin{solution}{exo:b3:complex-spectra-magnetism:12}
$T_{1/2} = 15\,000/75 = \qty{200}{K}$。在 \qty{250}{K} 时：$\Delta G = 15\,000 - 250 \times 75 = \qty{-3750}{J/mol}$，$K = \eu^{1.80} = 6.07$，
$x_{\mathrm{HS}} = 0.86$；$\chi_{\mathrm m}T = 0.86 \times 3.00 = \qty{2.6}{cm^3.K.mol^{-1}}$。
\end{solution}

\begin{solution}{pb:b3:complex-spectra-magnetism:1}
\textbf{1.} [Ar]$3d^3$.
\textbf{2.} $^4$F.
\textbf{3.} $^4$A$_{2g}$ + $^4$T$_{2g}$ + $^4$T$_{1g}$.
\textbf{4.} $^4$A$_{2g}$ 为 $t_{2g}^3$ 的谱项，三个电子全部位于较低的轨道中。
\textbf{5.} $^4$T$_{1g}$(P).
\textbf{6.} $^4$A$_{2g}\to{}^4$T$_{2g}$, $\to{}^4$T$_{1g}$(F), $\to{}^4$T$_{1g}$(P).
\textbf{7.} 561 和 \qty{414}{nm}。
\textbf{8.} 597 和 \qty{434}{nm}。
\textbf{9.} 红宝石吸收黄绿光和紫光，透过红光（略带
蓝光）；祖母绿的吸收带向红端移动，同样吸收橙红光，而绿光
透过。
\textbf{10.} 红宝石 \qty{17820}{cm^{-1}}，祖母绿 \qty{16740}{cm^{-1}}。
\textbf{11.} $^4$T$_{2g}$（$t_{2g}^2e_g$）与 $^4$A$_{2g}$（$t_{2g}^3$）具有相同的电子排斥能；它们的
差值即单电子跃迁能 $\Delta_{\mathrm o}$。
\textbf{12.} $B = (14\,305 - 547)/15 = \qty{917}{cm^{-1}}$.
\textbf{13.} \qty{614}{cm^{-1}}.
\textbf{14.} \qty{618}{cm^{-1}}.
\textbf{15.} 两者均为 0.67。
\textbf{16.} 29.0 和 27.1。
\textbf{17.} 红宝石的 $\nu_3 = \qty{38500}{cm^{-1}}$（\qty{260}{nm}），位于紫外区，该处的电荷转移
吸收掩盖了它。
\textbf{18.} 绿柱石的八面体间隙较大：较长的 Cr--O 键使分裂较小，
且分裂随距离急剧变化。
\textbf{19.} $10^7/695 = \qty{14390}{cm^{-1}}$, $23.4B$.
\textbf{20.} $^2$E$_g$ 主要取决于 $C$：测得能级更高，意味着 $C/B$ 更大，
同样用对角化可得约 5.3。
\textbf{21.} $^2$E$_g$ 与 $^4$A$_{2g}$ 同属 $t_{2g}^3$，仅相差一次自旋翻转：成键情况，
因而几何构型两者相同，该跃迁没有振动
序列（0--0 线）。
\textbf{22.} 该跃迁为自旋禁阻。
\textbf{23.} 其能量取决于 $B$ 和 $C$，几乎与 $\Delta_{\mathrm o}$ 无关（图中为一条水平线），
且两种宝石中的 $B$ 几乎相同。
\textbf{24.} 泵浦到宽吸收带中的离子经弛豫后，聚集在寿命较长的
$^2$E$_g$ 能级，使得其中的离子数可多于基态，即
粒子数反转。
\textbf{25.} \textbf{$\Delta_{\mathrm o}(\text{红宝石}) - \Delta_{\mathrm o}(\text{祖母绿}) = \qty{1080}{cm^{-1}}$}：
即红色与绿色之间的全部差异。
\end{solution}
